
\documentclass{article}
\usepackage{colortbl}
\usepackage{makecell}
\usepackage{multirow}
\usepackage{supertabular}

\begin{document}

\newcounter{utterance}

\centering \large Interaction Transcript for game `eqbench', experiment `eq\_bench', episode 127 with reflexion{-}Qwen3.5{-}9B.
\vspace{24pt}

{ \footnotesize  \setcounter{utterance}{1}
\setlength{\tabcolsep}{0pt}
\begin{supertabular}{c@{$\;$}|p{.15\linewidth}@{}p{.15\linewidth}p{.15\linewidth}p{.15\linewidth}p{.15\linewidth}p{.15\linewidth}}
    \# & \multicolumn{2}{c}{Player} && \multicolumn{2}{c}{Game Master} \\
    \hline

    \theutterance \stepcounter{utterance}  
    & & & \multicolumn{4}{p{0.6\linewidth}}{
        \cellcolor[rgb]{0.9,0.9,0.9}{
            \makecell[{{p{\linewidth}}}]{
                \texttt{\tiny{[P1$\langle$GM]}}
                \texttt{You have attempted this same task before and did not succeed. These are the plans you wrote after those attempts.} \\
\texttt{Trial \#1: New plan: I failed because I assigned a high score to "Dismissed" and "Angry" simultaneously, likely misinterpreting the nuance of Argyle's final outburst. Argyle's statement about not wanting to live in a "sterile, controlled environment" is a direct rebuttal to Serepta's accusation of being "carefree," indicating he feels his lifestyle and values are being invalidated and mocked, which aligns more strongly with feeling "Dismissed" (his perspective is ignored/dismissed) than purely "Angry" (though anger is present, the core feeling is the dismissal of his identity). Furthermore, "Proud" is incorrect because Argyle is defensive and resentful, not celebrating an achievement. "Inspired" is clearly 0. Next time, I will re{-}evaluate the specific emotional weight of the final exchange: Argyle feels his dreams are being belittled by Serepta's pragmatism, making "Dismissed" the dominant emotion (score 9{-}10), while "Angry" should be lower (score 3{-}4) as it is a reaction to the dismissal rather than the primary state. I will ensure the scores reflect this hierarchy and that the sum of non{-}zero scores feels proportionate to the intensity of the conflict.} \\
\\ 
\texttt{Trial \#2: New plan: I failed because I assigned a high score to "Dismissed" (9) while simultaneously giving a significant score to "Angry" (4), creating a conflicting emotional profile where the character feels both deeply ignored and moderately angry, rather than one dominant reaction to the other. Argyle's final outburst ("not everyone wants to live in a sterile, controlled environment") is a defensive rebuttal to Serepta's dismissal of his dreams, making "Dismissed" the primary driver of his state. The anger is a secondary, reactive component, not a parallel high{-}intensity emotion. To correct this, I must lower the "Angry" score to reflect its subordinate role (e.g., 2{-}3) and ensure "Dismissed" remains the highest score (9{-}10), while keeping "Inspired" and "Proud" at 0, as the dialogue is purely conflict{-}driven with no elements of inspiration or pride.} \\
\\ 
\\ 
\\ 
\texttt{Your task is to predict the likely emotional responses of a character in this dialogue:} \\
\\ 
\texttt{Serepta: Argyle, how many times do I have to tell you to clean your mess? } \\
\texttt{Argyle: Serepta, it's just a few sketches scattered around. It's not like I'm hoarding garbage.} \\
\texttt{Serepta: It's not about the sketches, Argyle. It's about respect. This is our shared space.} \\
\texttt{Argyle: Respect, huh? Like the respect you show when you rearrange my things without asking?} \\
\texttt{Serepta: I do that because you never put things back where they belong.} \\
\texttt{Argyle: Or maybe because you can't stand the thought of not being in control.} \\
\texttt{Serepta: That's not fair, Argyle. You know I have a lot on my plate.} \\
\texttt{Argyle: And you think I don't? } \\
\texttt{Serepta: You? You're always so carefree, always with that sketchbook in your hand. What do you know about pressure?} \\
\texttt{Argyle: You think I don't feel pressure? Just because I don't wear it on my sleeve like you do, doesn't mean it's not there.} \\
\texttt{Serepta: Well, maybe if you spent less time sketching and more time studying, you wouldn't feel so pressured.} \\
\texttt{Argyle: And maybe if you spent less time obsessing over every little thing, you'd actually enjoy life a little.} \\
\texttt{Serepta: I don't have the luxury of being carefree, Argyle. Some of us have responsibilities.} \\
\texttt{Argyle: And some of us have dreams, Serepta. But I guess you wouldn't understand that.} \\
\texttt{Serepta: I understand plenty, Argyle. More than you know.} \\
\texttt{Argyle: Yeah? Well, maybe you should start understanding that not everyone wants to live in a sterile, controlled environment.} \\
\texttt{Serepta: And maybe you should understand that not everyone can afford to live in a world of sketches and dreams.} \\
\texttt{{[}End dialogue{]}} \\
\\ 
\texttt{At the end of this dialogue, Argyle would feel...} \\
\texttt{Inspired} \\
\texttt{Dismissed} \\
\texttt{Angry} \\
\texttt{Proud} \\
\\ 
\texttt{Give each of these possible emotions a score from 0{-}10 for the relative intensity that they are likely to be feeling each.} \\
\\ 
\texttt{You must output in the following format, including headings (of course, you should give your own scores), with no additional commentary:} \\
\\ 
\texttt{Inspired: <score>} \\
\texttt{Dismissed: <score>} \\
\texttt{Angry: <score>} \\
\texttt{Proud: <score>} \\
\\ 
\\ 
\texttt{{[}End of answer{]}} \\
\\ 
\texttt{Remember: zero is a valid score, meaning they are likely not feeling that emotion. You must score at least one emotion > 0.} \\
\\ 
\texttt{Your answer:} \\
            }
        }
    }
     \\ \\

    \theutterance \stepcounter{utterance}  
    & \multicolumn{4}{p{0.6\linewidth}}{
        \cellcolor[rgb]{0.561,0.737,0.561}{
            \makecell[{{p{\linewidth}}}]{
                \texttt{\tiny{[P1$\rangle$GM]}}
                \texttt{Inspired: 0} \\
\texttt{Dismissed: 9} \\
\texttt{Angry: 3} \\
\texttt{Proud: 0} \\
            }
        }
    }
    & & \\ \\

    \theutterance \stepcounter{utterance}  
    & & & \multicolumn{2}{p{0.3\linewidth}}{
        \cellcolor[rgb]{0.9,0.9,0.9}{
            \makecell[{{p{\linewidth}}}]{
                \texttt{\tiny{[GM$|$GM]}}
                \texttt{Inspired: 0} \\
\texttt{Dismissed: 9} \\
\texttt{Angry: 3} \\
\texttt{Proud: 0} \\
            }
        }
    }
    & & \\ \\

    \theutterance \stepcounter{utterance}  
    & & & \multicolumn{2}{p{0.3\linewidth}}{
        \cellcolor[rgb]{0.9,0.9,0.9}{
            \makecell[{{p{\linewidth}}}]{
                \texttt{\tiny{[GM$|$GM]}}
                \texttt{Inspired: 0} \\
\texttt{Dismissed: 7} \\
\texttt{Angry: 8} \\
\texttt{Proud: 3} \\
            }
        }
    }
    & & \\ \\

    \theutterance \stepcounter{utterance}  
    & & & \multicolumn{2}{p{0.3\linewidth}}{
        \cellcolor[rgb]{0.9,0.9,0.9}{
            \makecell[{{p{\linewidth}}}]{
                \texttt{\tiny{[GM$|$GM]}}
                \texttt{game\_result = LOSE} \\
            }
        }
    }
    & & \\ \\

    \theutterance \stepcounter{utterance}  
    & & \multicolumn{4}{p{0.6\linewidth}}{
        \cellcolor[rgb]{0.95,0.95,0.95}{
            \makecell[{{p{\linewidth}}}]{
                \texttt{\tiny{[P1$\leftrightarrow$P]}}
                \texttt{Post{-}episode Reflexion (not sent to GM):} \\
\texttt{New plan: I failed because I assigned a high score to "Dismissed" (9) while simultaneously giving a significant score to "Angry" (3), creating a conflicting emotional profile where the character feels both deeply ignored and moderately angry, rather than one dominant reaction to the other. Argyle's final outburst ("not everyone wants to live in a sterile, controlled environment") is a defensive rebuttal to Serepta's dismissal of his dreams, making "Dismissed" the primary driver of his state. The anger is a secondary, reactive component, not a parallel high{-}intensity emotion. To correct this, I must lower the "Angry" score to reflect its subordinate role (e.g., 2{-}3) and ensure "Dismissed" remains the highest score (9{-}10), while keeping "Inspired" and "Proud" at 0, as the dialogue is purely conflict{-}driven with no elements of inspiration or pride.} \\
            }
        }
    }
    & \\ \\

\end{supertabular}
}

\end{document}
